\chapter{Convexity Adjustments and Constant-Maturity Products}\label{ch:rc:convexity-adjustments-and-constant-maturity-products}

A structured note pays, once a year, the ten-year euro swap rate observed that
day. For the coupon ten years from now the forward ten-year swap rate on
chapter~2's curves is 3.05\%, and a trader who quoted the coupon at the
forward would lose money: its fair value is 3.37\%. Nothing is wrong with the
forward. It is the fair rate of a swap starting in ten years, an expectation
under the measure whose numeraire is that swap's annuity; the note pays the
rate once, a year after it is observed, and under the measure of that
payment date the rate's expectation is higher, by an amount that depends on
the volatility of the rate and on how the annuity moves with it. The
difference is a convexity adjustment. Book 2 met one on futures (One Quant
Book 2, chapter 8); this chapter computes them for rates paid at the wrong
time, in the wrong currency or with the wrong maturity, prices them by
replication with the swaptions of chapter~5's cube, and values a note on the
difference of two swap rates.

\section{Paying a rate at the wrong time: timing adjustments}

A forward rate $F_k$ for $[T_{k-1},T_k]$ is a martingale under the
$T_k$-forward measure (\cref{prop:rc:vanilla-rates-options:measures}), whose
numeraire matches its natural payment date. Pay it at another date and its
expectation changes.

\begin{definition}[Timing adjustment]\label{def:rc:convexity-adjustments-and-constant-maturity-products:timing}
The \emph{timing adjustment}\index{timing adjustment} of a rate paid at a date
$T_p$ other than its natural payment date is
$\E^{T_p}[F_k(T_{k-1})]-F_k(0)$: the difference between its expectation under
the $T_p$-forward measure and its forward.
\end{definition}

\begin{proposition}[A rate paid when it fixes]\label{prop:rc:convexity-adjustments-and-constant-maturity-products:timing}
If $F_k$ is paid at its fixing date $T_{k-1}$ instead of $T_k$,
\[
\E^{T_{k-1}}\bigl[F_k(T_{k-1})\bigr] = F_k(0) +
\frac{\delta_k\,\Var^{T_k}\bigl(F_k(T_{k-1})\bigr)}{1+\delta_kF_k(0)} ,
\]
which is $F_k(0)+\delta_k\sigma_N^2T_{k-1}/(1+\delta_kF_k(0))$ in the normal model,
and $F_k(0)+\delta_kF_k(0)^2\sigma^2T_{k-1}/(1+\delta_kF_k(0))$ to first order in
the lognormal model.
\end{proposition}
\begin{proof}
The density of $\mathbb Q^{T_{k-1}}$ against $\mathbb Q^{T_k}$ on $\mathcal
F_{T_{k-1}}$ is the ratio of the numeraires normalised at zero,
$\frac{P(T_{k-1},T_{k-1})/P(T_{k-1},T_k)}{P(0,T_{k-1})/P(0,T_k)} =
\frac{1+\delta_kF_k(T_{k-1})}{1+\delta_kF_k(0)}$ (One Quant Book 4, chapter 5).
So $\E^{T_{k-1}}[F_k] = \E^{T_k}[F_k(1+\delta_kF_k)]/(1+\delta_kF_k(0))$, and
$\E^{T_k}[F_k^2] = F_k(0)^2+\Var^{T_k}(F_k)$.
\end{proof}

\begin{omfigure}
\begin{tikzpicture}[x=1.6cm,font=\small]
\draw[->,thick] (0,0) -- (6.4,0) node[right] {time};
\foreach \x/\l in {0/0,3/$T_{k-1}$,5/$T_k$} {\draw (\x,-0.1) -- (\x,0.1); \node[below] at (\x,-0.12) {\l};}
\draw[omDef,thick,fill=omDef!12] (3,0.45) rectangle (5,0.85); \node[font=\footnotesize] at (4,0.65) {accrual of $F_k$};
\draw[-{Stealth[length=2mm]},omDef,thick] (5,1.6) -- (5,0.95);
\node[text=omDef,align=center,font=\footnotesize] at (5,2.0) {natural payment:\\ $\E^{T_k}[F_k]=F_k(0)$};
\draw[-{Stealth[length=2mm]},omThm,thick] (3,1.6) -- (3,0.95);
\node[text=omThm,align=center,font=\footnotesize] at (3,2.0) {paid when fixed:\\ a timing adjustment};
\draw[-{Stealth[length=2mm]},omProp,thick] (5.9,1.6) -- (5.9,0.1);
\node[text=omProp,align=left,font=\footnotesize,anchor=west] at (6.0,1.3) {a swap rate paid once,\\ later: a CMS adjustment};
\end{tikzpicture}

\omcaption{Where a rate is paid decides the measure of its expectation. Paid
at the end of its accrual period, a forward rate needs no adjustment; paid at
its fixing, it needs a \omterm{def:rc:convexity-adjustments-and-constant-maturity-products:timing}{timing adjustment}; a swap rate paid once, on a single
date, needs a constant-maturity adjustment.}\label{fig:rc:convexity-adjustments-and-constant-maturity-products:timing}
\end{omfigure}

\begin{example}[A six-month rate paid in advance]\label{ex:rc:convexity-adjustments-and-constant-maturity-products:timing}
A six-month rate with a forward of 3\% fixing in five years, with a normal
volatility of 75 basis points (a \omterm{def:rc:vanilla-rates-options:lognormal}{lognormal volatility} of 25\%), paid at its
fixing: $0.5\times0.0075^2\times5/(1+0.5\times0.03) = 1.39$ basis points. Small
for one period, it is paid on every period of an in-arrears swap and grows
with the fixing date.
\end{example}

\section{Constant-maturity swaps and their replication by swaptions}

\begin{definition}[Constant-maturity swap]\label{def:rc:convexity-adjustments-and-constant-maturity-products:cms}
A \emph{constant-maturity swap}\index{constant-maturity swap} (CMS) exchanges
a floating leg whose rate at each reset is the par swap rate of a fixed tenor
(the ``CMS rate'', for instance the ten-year swap rate observed at the reset)
against a fixed or another floating leg. A CMS caplet pays
$\max(S(T)-K,0)$ on the CMS rate observed at $T$, once, at a date $T_p$.
\end{definition}

\begin{definition}[CMS convexity adjustment]\label{def:rc:convexity-adjustments-and-constant-maturity-products:cmsadj}
The \emph{CMS convexity adjustment}\index{CMS convexity adjustment} is
$\E^{T_p}[S_{a,b}(T)]-S_{a,b}(0)$: the expectation of the swap rate under the
forward measure of its payment date minus its forward, the expectation under
the annuity measure.
\end{definition}

Change measure from the annuity to the payment date:
\[
\E^{T_p}[S(T)] = \frac{A(0)}{P(0,T_p)}\,\E^{a,b}\Bigl[S(T)\,\frac{P(T,T_p)}{A(T)}\Bigr].
\]
The ratio $P(T,T_p)/A(T)$ is a function of the whole curve at $T$; to price, it
is mapped to a function of the swap rate alone.

\begin{definition}[Linear terminal swap-rate model]\label{def:rc:convexity-adjustments-and-constant-maturity-products:tsr}
The \emph{linear terminal swap-rate model}\index{linear terminal swap-rate model}
sets $P(T,T_p)/A(T) = a_0 + a_1S(T)$, with $a_0 = 1/\sum_i\delta_i$ (the value when
all rates are zero) and $a_1$ chosen so that the relation holds today:
$a_0+a_1S(0) = P(0,T_p)/A(0)$.
\end{definition}

\begin{proposition}[The CMS rate by replication]\label{prop:rc:convexity-adjustments-and-constant-maturity-products:cms}
Under the \omterm{def:rc:convexity-adjustments-and-constant-maturity-products:tsr}{linear terminal swap-rate model},
\[
\E^{T_p}[S(T)] = S(0) + \frac{a_1\,\Var^{a,b}(S(T))}{a_0+a_1S(0)},\qquad
\Var^{a,b}(S(T)) = 2\int_{-\infty}^{\infty} O(K)\,dK,
\]
where $O(K)$ is the out-of-the-money swaption price per unit annuity at strike
$K$ (a receiver below the forward, a payer above). With a flat normal smile,
$\Var^{a,b}(S(T)) = \sigma_N^2T$ (Hagan, 2003).
\end{proposition}
\begin{proof}
$\E^{a,b}[S(a_0+a_1S)] = a_0S(0) + a_1\bigl(S(0)^2+\Var^{a,b}(S)\bigr)$ since $S$ is
an $\mathbb Q^{a,b}$-martingale; divide by $a_0+a_1S(0) = P(0,T_p)/A(0)$. For the
variance, $S^2 = S(0)^2 + 2S(0)(S-S(0)) + 2\int_{-\infty}^{S(0)}(K-S)^+dK +
2\int_{S(0)}^{\infty}(S-K)^+dK$ (the static replication of One Quant Book 5,
chapter 14), and take expectations.
\end{proof}

\omcode{code/firm/cms/firm_cms.py}{63}{79}{The CMS rate by replication: out-of-the-money swaptions integrated over the smile.}\label{lst:rc:convexity-adjustments-and-constant-maturity-products:cms}

The CMS adjustment is therefore a price: a strip of out-of-the-money swaptions,
weighted by $2a_1$. It is paid for by the smile. With the cube's smile the
wings are richer than a \omterm{def:rc:vanilla-rates-options:flat}{flat volatility}, the variance is higher, and so is the
adjustment (\cref{fig:rc:convexity-adjustments-and-constant-maturity-products:adj}).

\begin{example}[Ten-year CMS on the euro cube]\label{ex:rc:convexity-adjustments-and-constant-maturity-products:cmsadj}
For the ten-year swap rate observed in one, two, three, five, seven and ten
years and paid one year later, the adjustments on chapter~2's curves are
1.39, 3.49, 5.87, 11.73, 17.40 and 26.58 basis points with each section's
at-the-money volatility held flat, and 1.62, 4.46, 7.92, 15.51, 22.91 and 32.14
with chapter~5's smile. The ten-year CMS rate observed in ten years has a
forward of 3.051\% and an expectation of 3.372\%: the hook's coupon.
\end{example}

\begin{omfigure}
\begin{tikzpicture}
\begin{axis}[width=11.5cm,height=5.4cm,
  xlabel={observation date of the ten-year swap rate (years)},ylabel={CMS adjustment (bp)},
  tick label style={font=\small},label style={font=\small},
  xmin=0,xmax=10.5,ymin=0,ymax=35,grid=major,grid style={gray!25},
  legend columns=2,legend style={font=\footnotesize,at={(0.5,-0.25)},anchor=north,draw=none,fill=none,/tikz/every even column/.append style={column sep=8pt}},
  table/col sep=comma]
\addplot[omDef,thick,mark=*] table[x=e,y=flat]{figdata/rates-credit-risk/06-convexity-adjustments-and-constant-maturity-products/cmsadj.csv};
\addplot[omThm,very thick,mark=square*] table[x=e,y=smile]{figdata/rates-credit-risk/06-convexity-adjustments-and-constant-maturity-products/cmsadj.csv};
\legend{flat at-the-money normal volatility,the cube's SABR smile}
\end{axis}
\end{tikzpicture}

\omcaption{Convexity adjustment of the ten-year euro CMS rate paid one year
after observation, by replication with a \omterm{def:rc:vanilla-rates-options:flat}{flat volatility} and with the smile of
chapter~5's cube. The smile adds 16 to 35 per cent to the adjustment. Data:
chapter~2's curves and chapter~5's synthetic cube; the chapter's
tutorial.}\label{fig:rc:convexity-adjustments-and-constant-maturity-products:adj}
\end{omfigure}

\begin{omfigure}
\begin{tikzpicture}
\begin{axis}[width=11.5cm,height=5.2cm,
  xlabel={swaption strike (\%)},ylabel={bp of adjustment per 1\% of strike},
  tick label style={font=\small},label style={font=\small},
  xmin=-3,xmax=7,ymin=0,ymax=9,grid=major,grid style={gray!25},
  table/col sep=comma]
\addplot[omDef,very thick] table[x=k,y=c]{figdata/rates-credit-risk/06-convexity-adjustments-and-constant-maturity-products/contrib.csv};
\draw[gray,dashed] (axis cs:3.051,0) -- (axis cs:3.051,9) node[pos=0.9,right,font=\footnotesize,text=black] {forward 3.05\%};
\end{axis}
\end{tikzpicture}

\omcaption{Where the ten-year CMS adjustment of the ten-year rate comes from:
the contribution of swaptions of each strike, receivers below the forward and
payers above, to the 32.14 basis points. Swaptions a percentage point or more
away from the money carry a large share of it, which is why the wings of the
smile matter. Data: the chapter's tutorial.}\label{fig:rc:convexity-adjustments-and-constant-maturity-products:contrib}
\end{omfigure}

A CMS caplet is replicated the same way. Writing $(S-K)^+(a_0+a_1S) =
(a_0+a_1K)(S-K)^+ + a_1((S-K)^+)^2$ and $((S-K)^+)^2 = 2\int_K^\infty(S-L)^+dL$,
its value is $A(0)\bigl[(a_0+a_1K)\,\mathrm{Pay}(K) + 2a_1\int_K^\infty
\mathrm{Pay}(L)\,dL\bigr]$, with $\mathrm{Pay}$ the payer price per unit annuity.

\omcode{code/firm/cms/firm_cms.py}{82}{90}{A CMS caplet as payer swaptions at its strike and above.}\label{lst:rc:convexity-adjustments-and-constant-maturity-products:caplet}

\begin{example}[A CMS caplet]\label{ex:rc:convexity-adjustments-and-constant-maturity-products:caplet}
A caplet at 3.50\% on the ten-year rate observed in five years (forward
3.098\%, smile volatility 75.9 basis points at the strike), paid a year later,
is worth 48.5 basis points of notional by replication; discounting a
Bachelier call on the forward would give 43.0, 13\% too little.
\end{example}

\section{The linear terminal swap-rate model in practice}

The model needs only today's curve: $a_0$ from the accruals, $a_1$ from
$P(0,T_p)/A(0)$. It is exact when the curve at $T$ is flat and moves in parallel
and is a good approximation otherwise; its alternatives map $P(T,T_p)/A(T)$
through a flat-curve annuity formula (Hagan's standard model) or through a
short-rate model's bonds (chapter~7). What it misses is the dependence of the
mapping on the curve's shape, which matters for long payment delays and steep
curves; what matters more in practice is the smile far from the money, where
replication reaches. A desk that marks CMS products therefore extrapolates its
cube's wings with care (chapter~5's densities) and often calibrates the wing
parameters to traded CMS spreads rather than to swaptions.

\section{Quanto adjustments in rates}

A euro swap rate paid in dollars is a rate observed in one currency and paid in
another. Under the dollar payment measure its drift changes by the covariance
with the exchange rate, the quanto adjustment of One Quant Book 5, chapter 17.
For a normal rate with volatility $\sigma_N$, an EURUSD volatility $\sigma_X$
(dollars per euro) and correlation $\rho$ between the rate and EURUSD, the
expectation moves by $-\rho\,\sigma_N\sigma_X\,T$ to first order.

\begin{example}[A euro rate paid in dollars]\label{ex:rc:convexity-adjustments-and-constant-maturity-products:quanto}
With $\sigma_N = 75$ basis points, $\sigma_X = 8\%$, $\rho = 0.3$ and five years,
the adjustment is $-0.3\times0.0075\times0.08\times5 = -9$ basis points: when euro
rates rise the euro tends to rise too, so each basis point is paid in more
valuable euros than the dollar payment delivers.
\end{example}

\section{CMS spread options}

\begin{definition}[CMS spread option]\label{def:rc:convexity-adjustments-and-constant-maturity-products:spread}
A \emph{CMS spread option}\index{CMS spread option} pays $\max(S_1(T)-S_2(T)-K,0)$
at $T_p$ on the difference of two CMS rates of different tenors observed at the
same date, typically the ten-year and the two-year; a strip of them with
$K=0$ pays a steepener coupon.
\end{definition}

Priced with each rate normal around its own CMS expectation, the spread is
normal with volatility $\sigma_s = \sqrt{\sigma_1^2+\sigma_2^2-2\rho\sigma_1\sigma_2}$:
the price depends on the correlation between the two rates, which no swaption
quotes. The desk that sells a steepener note is short spread volatility and so
long correlation; its hedge is correlation-sensitive products, and its mark
depends on an estimate.

\begin{example}[The spread's volatility]\label{ex:rc:convexity-adjustments-and-constant-maturity-products:spreadvol}
Over 2016--2026 daily changes of the two- and ten-year Treasury par yields had
a correlation of 0.769 (chapter~3's data, a proxy for swap rates). With the
cube's at-the-money volatilities for the one-year expiry, the spread of the
ten- and two-year euro rates has a normal volatility of 45.1 basis points, well
below either rate's.
\end{example}

\begin{omfigure}
\begin{tikzpicture}
\begin{axis}[width=11cm,height=5.2cm,
  xlabel={correlation of the ten- and two-year rates},ylabel={fair participation},
  tick label style={font=\small},label style={font=\small},
  xmin=0.48,xmax=1.0,ymin=2,ymax=5.4,grid=major,grid style={gray!25},
  table/col sep=comma]
\addplot[omThm,very thick,mark=*,mark size=1.5pt] table[x=rho,y=p]{figdata/rates-credit-risk/06-convexity-adjustments-and-constant-maturity-products/steep.csv};
\draw[gray,dashed] (axis cs:0.769,2) -- (axis cs:0.769,5.4) node[pos=0.9,left,font=\footnotesize,text=black] {Treasury estimate 0.769};
\end{axis}
\end{tikzpicture}

\omcaption{Fair participation of a ten-year steepener note (the weekend
problem) as a function of the correlation between the ten- and two-year swap
rates: the higher the correlation, the less volatile the spread and the cheaper
each unit of its floor, so the more of it the note can pay. Data: the chapter's
tutorial.}\label{fig:rc:convexity-adjustments-and-constant-maturity-products:steep}
\end{omfigure}

\section{Tutorial: replicating a CMS rate}\label{tut:rc:convexity-adjustments-and-constant-maturity-products:tutorial}

\begin{tutorial}
\textbf{Goal.} Compute CMS adjustments by replication on the cube, price a CMS
caplet and a steepener note. \textbf{End state:}
\cref{fig:rc:convexity-adjustments-and-constant-maturity-products:adj,fig:rc:convexity-adjustments-and-constant-maturity-products:contrib,fig:rc:convexity-adjustments-and-constant-maturity-products:steep}.
\begin{tutsteps}
\item \textbf{Set-up}: \texttt{setup(e, tenor)} returns the forward swap rate, its
annuity, the payment discount factor and the accruals from chapter~2's curves.
\item \textbf{Adjustments}: \texttt{cms\_table()} with a flat smile
(\texttt{cms\_rate\_hagan}) and with the cube (\texttt{cms\_rate}); the cube's
shifted model has no mass below $-3\%$, so the integral starts there.
\item \textbf{Where it comes from}: \texttt{replication\_contributions()}.
\item \textbf{The note}: \texttt{steepener()}; \texttt{fig\_rc\_cms.py} writes the
charts.
\end{tutsteps}
\textbf{What to change next.} Replace the cube's smile by its flat at-the-money
volatility in the steepener and compare; move the payment date to the end of
the underlying swap and watch the adjustment grow.
\end{tutorial}

\section{Build: CMS pricing}\label{bld:rc:convexity-adjustments-and-constant-maturity-products:cms}

\begin{build}
\bfield{Purpose} Pricing of every product paying a rate at the wrong time, in
the wrong currency or with a constant maturity; the CMS legs of the book's
structured notes.
\bfield{Interface} \texttt{timing\_adjustment\_black},
\texttt{timing\_adjustment\_normal}; \texttt{linear\_tsr(annuity0, df\_pay, fwd,
accruals)}; \texttt{otm\_integral}, \texttt{cms\_rate}, \texttt{cms\_caplet},
\texttt{cms\_rate\_hagan}; \texttt{quanto\_adjustment\_normal};
\texttt{spread\_vol}, \texttt{cms\_spread\_option}.
\bfield{Rules} The smile is any function from strike to normal volatility (the
cube of chapter~5 in the book); the integral is bounded below by the smile
model's floor; numpy only.
\bfield{Acceptance tests} \texttt{code/firm/cms/tests/}: with a flat smile the
replicated adjustment equals the closed form; with no slope in the annuity
mapping there is no adjustment; a deep in-the-money CMS caplet is the
discounted CMS forward; \omterm{def:rc:convexity-adjustments-and-constant-maturity-products:timing}{timing adjustments} agree between the Black and normal
forms at matched volatilities; perfect correlation removes the spread option's
time value.
\bfield{Stretch} Hagan's standard (flat-curve annuity) mapping; \omterm{def:rc:convexity-adjustments-and-constant-maturity-products:spread}{CMS spread options}
with a copula of the two smiles; a CMS swap's full leg with payment delays.
\end{build}

\begin{omsources}
\item P.\ S.\ Hagan, ``Convexity conundrums: pricing CMS swaps, caps, and
floors'', \emph{Wilmott Magazine}, March 2003.
\item L.\ Andersen and V.\ Piterbarg, \emph{Interest Rate Modeling}, volume 3,
Atlantic Financial Press, 2010, chapters 16--17.
\item US Department of the Treasury, \emph{Daily Treasury Par Yield Curve Rates}
(the correlation estimate).
\end{omsources}

\section{Exercises}

\begin{exercise}[$\star$]\label{exo:rc:convexity-adjustments-and-constant-maturity-products:1}
A three-month rate of 2.5\% fixing in four years, with 80 basis points of normal
volatility, is paid at its fixing. Compute its \omterm{def:rc:convexity-adjustments-and-constant-maturity-products:timing}{timing adjustment}.
\end{exercise}

\begin{exercise}[$\star$]\label{exo:rc:convexity-adjustments-and-constant-maturity-products:2}
Why is the CMS adjustment positive for a rate paid shortly after it is observed?
What sign does $a_1$ have when rates are positive?
\end{exercise}

\begin{exercise}[$\star$]\label{exo:rc:convexity-adjustments-and-constant-maturity-products:3}
For a ten-year annual swap with accruals of 1.0139 each, compute $a_0$.
\end{exercise}

\begin{exercise}[$\star\star$]\label{exo:rc:convexity-adjustments-and-constant-maturity-products:4}
With a flat smile, the five-year adjustment of \cref{ex:rc:convexity-adjustments-and-constant-maturity-products:cmsadj}
is 11.73 basis points. How large would it be at a normal volatility 20\% higher?
\end{exercise}

\begin{exercise}[$\star\star$]\label{exo:rc:convexity-adjustments-and-constant-maturity-products:5}
Read \cref{fig:rc:convexity-adjustments-and-constant-maturity-products:contrib}:
roughly what share of the adjustment comes from swaptions more than two
percentage points from the forward? What does that say about how the cube's
wings are marked?
\end{exercise}

\begin{exercise}[$\star\star$]\label{exo:rc:convexity-adjustments-and-constant-maturity-products:6}
In \cref{ex:rc:convexity-adjustments-and-constant-maturity-products:quanto}, what
happens to the adjustment if the correlation is $-0.3$, and why?
\end{exercise}

\begin{exercise}[$\star\star\star$]\label{exo:rc:convexity-adjustments-and-constant-maturity-products:7}
\emph{Coding.} Compute the ten-year CMS adjustment of the rate observed in five
years when it is paid at the end of the underlying swap (fifteen years) instead
of one year after observation. Explain the sign.
\end{exercise}

\begin{exercise}[$\star\star\star$]\label{exo:rc:convexity-adjustments-and-constant-maturity-products:8}
\emph{Find the flaw.} ``A CMS swap is linear in rates, so we value its coupons at
the forward swap rates and hedge them with forward-starting swaps.'' Correct it.
\end{exercise}

\section{Problem: The Steepener Note}

\begin{problem}[{Weekend problem --- pricing a ten-year steepener}]\label{pb:rc:convexity-adjustments-and-constant-maturity-products:1}
A bank issues a ten-year euro note that pays each year $k=1,\dots,10$ a coupon of
$p\times\max(S_{10}(k)-S_2(k),0)$, the ten-year minus the two-year swap rate
observed at year $k$, paid at $k+1$. The bank's funding target is to pay coupons
worth as much as a fixed 2\% coupon on the same dates. Rates are those of
chapter~2, smiles those of chapter~5, and the correlation between the two rates
is estimated at 0.769.

\medskip\noindent\textbf{Part I --- The pieces.}
\begin{enumerate}
\item What is each coupon, as an option?
\item Why does each rate need a CMS adjustment, and which one is larger?
\item Give the value of the fixed 2\% leg per unit notional.
\item Give the value of the coupon leg with $p=1$.
\item Give the fair participation $p$.
\end{enumerate}

\medskip\noindent\textbf{Part II --- Correlation.}
\begin{enumerate}[resume]
\item Give the fair participation at correlations of 0.6, 0.8 and 0.95.
\item Why does it rise with correlation?
\item Which way is the issuing bank exposed to correlation after selling the note?
\item Give the normal volatility of the first year's spread.
\item Why is Treasury data only a proxy here?
\end{enumerate}

\medskip\noindent\textbf{Part III --- Risk.}
\begin{enumerate}[resume]
\item Which swaptions hedge the CMS adjustments?
\item What happens to the note's value if the smile's wings rise?
\item Is the note's value sensitive to the level of rates, the slope, or both?
\item What happens to the investor if the curve inverts?
\item Why do dealers' books of such notes become large and one-sided?
\end{enumerate}

\medskip\noindent\textbf{Part IV --- Judgement.}
\begin{enumerate}[resume]
\item What participation would you quote, and why not 2.96?
\item How would you mark correlation in the absence of quotes?
\item What would you tell an investor who says the curve ``always'' steepens back?
\item State the \emph{named result}: the fair participation at the estimated
correlation, and its range between correlations of 0.6 and 0.95.
\item In one sentence: what does a steepener investor buy?
\end{enumerate}
\end{problem}

\section{Interview questions}

\begin{interviewq}[$\star$ \iqroles{researcher, trader}]\label{iq:rc:convexity-adjustments-and-constant-maturity-products:1}
What is a convexity adjustment, in one sentence, and give two examples.
\end{interviewq}

\begin{interviewq}[$\star\star$ \iqroles{researcher}]\label{iq:rc:convexity-adjustments-and-constant-maturity-products:2}
Derive the \omterm{def:rc:convexity-adjustments-and-constant-maturity-products:timing}{timing adjustment} of a rate paid at its fixing date.
\end{interviewq}

\begin{interviewq}[$\star\star$ \iqroles{researcher, bank}]\label{iq:rc:convexity-adjustments-and-constant-maturity-products:3}
How do you price a CMS caplet from swaptions?
\end{interviewq}

\begin{interviewq}[$\star\star$ \iqroles{trader}]\label{iq:rc:convexity-adjustments-and-constant-maturity-products:4}
You are short CMS coupons. Which way do you want implied volatility to move, and
which strikes matter most?
\end{interviewq}

\begin{interviewq}[$\star\star\star$ \iqroles{researcher, risk}]\label{iq:rc:convexity-adjustments-and-constant-maturity-products:5}
A \omterm{def:rc:convexity-adjustments-and-constant-maturity-products:spread}{CMS spread option} depends on a correlation that is not quoted. How do you mark
and hedge it?
\end{interviewq}

\begin{interviewq}[$\star\star\star$ \iqroles{developer}]\label{iq:rc:convexity-adjustments-and-constant-maturity-products:6}
Your CMS replication gives a different number every time the strike grid
changes. What do you check?
\end{interviewq}
